\documentclass[a4paper]{article}

\usepackage{graphicx}

\usepackage{FillAFour}
\usepackage{HaskellVerb}

\begin{document}

\begin{titlepage}
    \begin{center}
        \vspace*{1cm}
            
        \Huge
        \textbf{JSON Parser in Haskell}
            
        \vspace{0.5cm}
        \LARGE
        Assignment 2
            
        \vspace{1.5cm}
            
        \textbf{Sam Grumley}
            
        \vfill
            
        Programming Languages\\
         3802ICT\\
            
        \vspace{0.8cm}
            
            
        \Large
         Griffith University\\
         School of Information and Communication Technology\\
         S5048240\\
         5/10/2020
            
    \end{center}
\end{titlepage}

\section{Introduction}
This report both documents and demonstrates the implementation of a JSON file parser writen in Haskell.
The EBNF and syntax diagrams have been constructed by the program Syntrax. 
The implementation is built around the third party module: ABR.

\section{Instructions}
Run the below commands in terminal to compile and parse your chosen JSON file
\begin{code}
ghc json-parser.hs 
json-parser input.json
\end{code}


\section{JSON EBNF}
EBNF and syntax diagrams:

\subsection{JSON Object}
\EBNFInput{syntax/json.ebnf}
{\centering
   \includegraphics[scale=0.9]{syntax/json}
}

\subsection{JSON Value}
\EBNFInput{syntax/value.ebnf}
{\centering
   \includegraphics[scale=0.9]{syntax/value}
}

\subsection{JSON Array}
\EBNFInput{syntax/array.ebnf}
{\centering
   \includegraphics[scale=0.9]{syntax/array}
}

\subsection{JSON Number}
JNumber was not manually written as the ABR library handles number with the funcion floatL.
\EBNFInput{syntax/number.ebnf}
{\centering
   \includegraphics[scale=0.9]{syntax/number}
}

\subsection{JSON String}
JString was not manually written as the ABR library handles number with the funcion stringL.
\EBNFInput{syntax/string.ebnf}
{\centering
   \includegraphics[scale=0.9]{syntax/string}
}

\subsection{JSON Bool}
JString was not manually written as the ABR library handles number with the funcion stringL.
\EBNFInput{syntax/bool.ebnf}
{\centering
   \includegraphics[scale=0.9]{syntax/bool}
}

\subsection{JSON Object Symbols}
JString was not manually written as the ABR library handles number with the funcion stringL.
\EBNFInput{syntax/objectSymbols.ebnf}
{\centering
   \includegraphics[scale=0.9]{syntax/objectSymbols}
}

\subsection{JSON Array Symbols}
JString was not manually written as the ABR library handles number with the funcion stringL.
\EBNFInput{syntax/arraySymbols.ebnf}
{\centering
   \includegraphics[scale=0.9]{syntax/arraySymbols}
}

\subsection{JSON Reapeatable Symbols}
JString was not manually written as the ABR library handles number with the funcion stringL.
\EBNFInput{syntax/repeatableSymbols.ebnf}
{\centering
   \includegraphics[scale=0.9]{syntax/repeatableSymbols}
}



\section{Implementation}

\subsection{Data Structure}
\input{data.lhs}

\subsection{Lexers}
\input{lexers.lhs}

\subsection{Parser}
\input{parsers.lhs}


\end{document}
